Electroencephalograms (EEGs) are recordings of the brain activity. Electrodes
are put on the skull to record electric activty in multiple places.
\cref{figeegsetup} shows such a device. EEGs are very widely used with machine
learning models to retrieve relevent information from the brain, together they
help to predict epilepsy seizures, diagnose Parkinson disease, or control remote
devices with mind \parencite{sahota2023}.

The field of generative AI is growing very rapidly in the recent years and its
usage for image as well as text is now very well known to the public.  This tool
rapidly proved very promising usage in the neuroscience as it could be used to
learn from our brain signals and produce images or texts, the big breakthrough
in image generation from \cite{diffusion}. However, DreamDiffusion isn't
semantically correct, during the internship it was attempted to propose a new
metric for EEG to image generative models to measure their semantic fidelity.

\insertfigure
    {assets/images/eeg-setup}
    {height=0.2\textheight}
    {An EEG cap}
    {figeegsetup}

\subsection*{The internship}

The internship aimed at evalutating the coherence and design metrics in EEG to
image reconstruction models and submit a research publication. It took place in
Hyderabad, Telangana, India, in Mahindra University (MU), under the supervision
of Dr. Nidhi Goyal. It was in a team of researchers made of interns, researchers
from both academics and industry, and one B.Tech volunteed.

Before the internship it was asked to complete the
\href{http://cs231n.stanford.edu}{CS231n} course "Deep Learning for Computer
Vision" from Stanford University, all the lecture videos were watched but only 2
assignements out of the 3 were completed. This course was given in 2017 at
Stanford University and covered many topics that went quite deeply into deep
learning, images, generation. The assignements were about implementing models
and training them from scratch.

During the internship multiple events were attended.
\begin{itemize}
\item
Being a teaching assistant for two practicals on Natural Language Processing
(NLP) during the \href
{https://www.mahindrauniversity.edu.in/events/aic-mahindra-mahindra-university-in-collaboration-with-the-all-india-council-for-technical-education-aicte-and-the-ministry-of-educations-innovation-cell-mic-successfully-conducted-the-facult/}
{Faculty Development Program}, the task was to go to the participants, solve the
issues they had and answer their questions.

\item
Attending the Generative AI for Image Processing and Healthcare
(GAIIH2025)\cite{gaiih2025} a 3 days summer school on Generative AI for
Healthcare, this completed. Many lectures explained quite deeply some machine
learning concepts, with hands-on sessions to implement them from scratch. The
topics it covered were very recent like Diffusion models, Stable Diffusion,
GANs, LLMs, Prompt Engineering, RNNs, Transformers, Visual Transformers,
AutoEncoders, U-Nets, LORA, Multimodal Models.

\item
A 2 days workshop by NVidia was attended, going through the different tools that
are useful for Machine Learning researchers. They showed how to use their DGX
cluster and run some models and pipelines.
\end{itemize}

For the internship, an access to the University's supercomputer, a NVidia
DGX-A100 cluster was given. Learning how to use it and paralellize work on the
two GPUs took some time.

One goal of the internship was to try to publish a research paper in a A/A*
conference (NeurIPS, WACV or AAAI). At the end of the internship the paper
wasn't mature enough to be published so two new interns in MU continued the
project, still with the help of Sai Ankit Reddy, Dr Nidhi Goyal and Dr Sai Amrit
Patnaik. We organized videoconferences to discuss the direction of the paper and
the experiments that had to be conducted.  It resulted in submitting the paper
in a conference where it is today still under review.

During the course of the internship there was also a side project mixing
optometry, neuroscience and data science I was involved in together with Gaspard
Lefort, lead by Dr Shiva Ram Reddy from University of Hyderabad (UoH) , Dr
Sunder Bukya and Dr Nidhi Goyal both from MU.  It was about studying distracted
driving behaviours among indian adults.